%% 2026-10-27 に ~/slides/common/ から写した（MER-DA を自己完結させるため）．以後は MER-DA/slides/common/ のこのファイルを正本とし，~/slides とは同期しない．
%% =====================================================================
%%  common/macros.tex — 全スライド共通のマクロ
%%  ここに置くのは「どの発表でも使うもの」だけ．
%%  その発表固有のマクロは deck 側の main.tex に書く．
%% =====================================================================

% ---------- 強調 ------------------------------------------------------
\newcommand{\key}[1]{\textcolor{mAccentBlue}{\textbf{#1}}}  % 提案手法・注目対象
\newcommand{\good}[1]{\textcolor{mGood}{#1}}                % 改善・良い結果
\newcommand{\bad}[1]{\textcolor{mBad}{#1}}                  % 劣化・悪い結果
\newcommand{\muted}[1]{\textcolor{mMuted}{#1}}              % 補足

% 出典表示（右下に小さく．スライドの引用はこれで足りることが多い）
\newcommand{\source}[1]{%
  \vspace{-0.3em}\hfill{\tiny\color{mMuted}[#1]}\par}

% 「ひとことメッセージ」帯：そのスライドで言いたいことを1行で置く
\newcommand{\takeaway}[1]{%
  \vspace{0.4em}%
  \begin{beamercolorbox}[wd=\textwidth,sep=0.5em,rounded=true]{block body}%
    \textbf{#1}%
  \end{beamercolorbox}}

% ---------- 数式（機械学習でよく使うもの）-----------------------------
\newcommand{\R}{\mathbb{R}}
\newcommand{\E}{\mathbb{E}}
\newcommand{\Prob}{\mathbb{P}}
\newcommand{\vect}[1]{\bm{#1}}
\newcommand{\mat}[1]{\bm{#1}}
\newcommand{\normtwo}[1]{\left\lVert #1 \right\rVert_2}
\newcommand{\norm}[2][]{\left\lVert #2 \right\rVert_{#1}}
\newcommand{\abs}[1]{\left\lvert #1 \right\rvert}
\newcommand{\inner}[2]{\left\langle #1,\, #2 \right\rangle}
\DeclareMathOperator*{\argmin}{arg\,min}
\DeclareMathOperator*{\argmax}{arg\,max}
\newcommand{\loss}{\mathcal{L}}
\newcommand{\dataset}{\mathcal{D}}
\newcommand{\model}{f_{\bm{\theta}}}
\newcommand{\given}{\,\vert\,}

% ---------- 図の貼り込み ----------------------------------------------
% \fig{ファイル名}{幅(\textwidth比)}  例: \fig{overview}{0.8}
\newcommand{\fig}[2]{%
  \begin{center}\includegraphics[width=#2\textwidth]{#1}\end{center}}

% \figc{ファイル名}{幅}{キャプション}
\newcommand{\figc}[3]{%
  \begin{figure}\centering
    \includegraphics[width=#2\textwidth]{#1}
    \caption{#3}
  \end{figure}}

% ---------- 表 --------------------------------------------------------
% 数値表の最良値を強調する（\best{0.812}）
\newcommand{\best}[1]{\textbf{\textcolor{mAccentBlue}{#1}}}

% ---------- 進捗報告でよく使う印 --------------------------------------
\newcommand{\done}{\textcolor{mGood}{\textbf{[完了]}}}
\newcommand{\wip}{\textcolor{mLightBrown}{\textbf{[進行中]}}}
\newcommand{\todo}{\textcolor{mMuted}{\textbf{[未着手]}}}
\newcommand{\blocked}{\textcolor{mBad}{\textbf{[停滞]}}}
